\Needspace{8\baselineskip}
\section{The increment relative to the inspected repository}
\label{sec:delta}

\subsection{What is already present, and therefore not a contribution here}

The reviewed Forge snapshot is a consolidated design and prototype collection,
not an installed Lean tactic. Its article already covers conclusion-directed
retrieval, structural induction and accumulator generalization, auxiliary-lemma
search, polynomial recurrence invariants, affine and integer witnesses,
nonlinear certificates, and Boolean/theory cooperation. It also incorporates a
substantial review of Leant and Djex, including their verification boundaries
and behavioral-evidence distinctions~\cite{forge,forge-structural,forge-provenance,forge-neighbors}.

Consequently, ``add induction'', ``use a computer algebra system'', ``produce
certificates'', and ``reuse Djex'' would not answer the present engineering
question. This report instead supplies three concrete workers and explains the
mathematics and interfaces required to turn their finite outputs into proofs of
unbounded behavior.

\begin{longtable}{@{}p{0.23\textwidth}p{0.32\textwidth}p{0.39\textwidth}@{}}
\toprule
Area & Reviewed Forge design & Increment developed here \\
\midrule\endhead
Polynomial invariants & Polynomial additive recurrences and conserved quantities;
general ideal multipliers mentioned as a next step & An executable descending
ideal-preimage algorithm, tracked multiplier extraction, and a greatest-fixed-point
theorem for a specified bounded generator class. \\
Recursive equivalence & Induction, rewriting, and lemma discovery & A finite
reachable-space decision procedure for rational weighted representations,
including tensor-composed observations and checkable counterexample words. \\
Finite sums & Polynomial finite-difference invariants & A concrete
creative-telescoping ansatz for hypergeometric binomial-power sums, including
moving support, pole regularization, boundary checks, and recurrence uniqueness. \\
Leant/Djex integration & Typed candidates and checked behavioral domains &
Proposed polynomial-state and weighted-observation domains tied to the existing
source-owned candidate graphs; a separate composition experiment for blocked
higher-order synthesis. \\
\bottomrule
\end{longtable}

This is an extension relative to the \emph{reviewed consolidated design}, not a
claim that no analogous phrase occurs in any historical proposal or that the
underlying mathematical subjects are new. Polynomial-ideal invariant generation,
weighted-automaton equivalence, and creative telescoping all have substantial
prior literature~\cite{sankaranarayanan,kiefer,brochet}. The contribution here is
the explicit selection of useful fragments, executable algorithms, finite
certificate formats, source-to-Lean obligations, and tests tailored to Forge.

\subsection{Pinned inputs and evidence levels}

Repository identities are recorded in \code{results/status.json}. The relevant
snapshots are:
\begin{center}
\begin{tabular}{@{}ll@{}}\toprule
Repository & Commit \\\midrule
Forge & \texttt{674521027d968d59f7b83220ed52304a85cb55e2} \\
Leant & \texttt{6bf05ad78c467989e68290f2d08bbed40802d485} \\
Djex & \texttt{e8778f4ebd63e1f9b9b410fa4de8d14a8a04c9e5} \\
\bottomrule\end{tabular}
\end{center}
The Forge README records Lean 4.34.0 and Mathlib
\texttt{1cf325a0cf67aca2b04d76b5380ff6a9e410aefa} as its baseline. It also records
three elaborated Lean-core files, which must not be confused with compiled
certificate workers~\cite{forge}. The experiments in this package used Python
3.13.5 and SymPy 1.14.0. Neither Lean nor Lake was available in the execution
runtime. The supplied Lean specimens therefore have status \code{NOT_RUN}.

Four kinds of statement occur in the rest of the report. A theorem accompanied
by a written proof is a mathematical claim about an explicitly defined
algorithm or certificate language. A recorded experiment is an observation
about the included Python programs. An integration design identifies code still
to write. A compiled Lean proof would be a fourth, stronger artifact; this
package does not contain one. This distinction is especially important when a
Python checker verifies a polynomial identity: the identity is useful proof
data, but the checker invocation is not itself a Lean theorem.

\subsection{A common target: finite closure of an infinite behavior}

The three workers share a useful mathematical pattern. Instead of expanding
an unbounded sequence of recursive calls, find a finite object that is closed
under the operation generating those calls.

For invariants, the object is an ideal generated by finitely many low-degree
polynomials. Applying a transition to a generator lands back in that ideal,
possibly modulo an equality guard. For a weighted fold, the object is a
finite-dimensional space of residual state vectors. Appending a letter keeps
the state in that space, and the output vanishes throughout it. For a definite
sum, a finite linear combination of parameter shifts is a difference in the
summation index. Summing cancels the interior and leaves explicit endpoints.

These are not three names for the same implementation. Their expensive steps
are different: Gr\"obner reduction, rational linear algebra, and polynomial
coefficient solving with support analysis. Their certificates are different,
too. What is shared is that the final induction or summation proof need not
reproduce the search.

\begin{finding}
The strongest integration point is not ``try a bigger solver on the original
goal''. It is ``discover a finite closed representation, prove its bridge to the
source expression, and discharge the resulting finite identities''.
\end{finding}

\subsection{Representative goals}

The ideal worker should prove that a state initialized on $y=x^2$ stays on that
curve under either $(x,y)\mapsto(x+1,y+2x+1)$ or
$(x,y)\mapsto(x^2,y^2)$. The second transition does not conserve the polynomial
$x^2-y$: it multiplies it by $x^2+y$. That distinction defeats a conservation-only
ansatz without making the desired invariant difficult.

The weighted worker should prove, for every binary word $w$,
\[
 \#ab(w)+\#ba(w)=\#a(w)\,\#b(w),
\]
where $\#ab$ and $\#ba$ count ordered two-letter subsequences, not adjacent
substrings. One side is a five-state affine fold; the other is a tensor product
of two counters. The proof comes from a rank-five closure in a nine-dimensional
difference representation.

The telescoping worker should discover and certify, for
$S(n)=\sum_{k=0}^{n}\binom nk^3$,
\[
 (n+2)^2S(n+2)-(7n^2+21n+16)S(n+1)-8(n+1)^2S(n)=0.
\]
It should not accept a rational summation certificate merely because its generic
formula simplifies: the moving endpoints intersect apparent poles. The explicit
regularization and boundary checks in Section~\ref{sec:telescoping} are part of
the capability, not optional cleanup.
